\SourceLeaf{245}%
wie im vorigen Artikel, sofern sie sich allgemein auf einen unbestimmten Punkt einer beliebigen kürzesten Linie beziehen.

Alle kürzesten Linien derselben Länge \(r\) enden auf einer anderen Linie. Ihre von einem beliebigen Anfangspunkt aus gemessene Länge bezeichnen wir mit \(v\). Daher lässt sich \(v\) als Funktion der Unbestimmten \(r,\varphi\) auffassen. Bezeichnen wir mit \(\lambda'\) den Punkt auf der Kugeloberfläche, der der Richtung des Elements \(\mathrm{d}v\) entspricht, und mit \(\xi',\eta',\zeta'\) seine auf den Mittelpunkt der Kugel bezogenen Koordinaten, so gilt:
\[
\frac{\mathrm{d}x}{\mathrm{d}\varphi}=\xi'\mathord{.}\frac{\mathrm{d}v}{\mathrm{d}\varphi},\quad
\frac{\mathrm{d}y}{\mathrm{d}\varphi}=\eta'\mathord{.}\frac{\mathrm{d}v}{\mathrm{d}\varphi},\quad
\frac{\mathrm{d}z}{\mathrm{d}\varphi}=\zeta'\mathord{.}\frac{\mathrm{d}v}{\mathrm{d}\varphi}
\]
Hieraus und aus
\[
\frac{\mathrm{d}x}{\mathrm{d}r}=\xi,\quad\frac{\mathrm{d}y}{\mathrm{d}r}=\eta,\quad\frac{\mathrm{d}z}{\mathrm{d}r}=\zeta
\]
folgt
\[
\begin{aligned}
\frac{\mathrm{d}x}{\mathrm{d}r}\mathord{.}\frac{\mathrm{d}x}{\mathrm{d}\varphi}+\frac{\mathrm{d}y}{\mathrm{d}r}\mathord{.}\frac{\mathrm{d}y}{\mathrm{d}\varphi}+\frac{\mathrm{d}z}{\mathrm{d}r}\mathord{.}\frac{\mathrm{d}z}{\mathrm{d}\varphi}
&=(\xi\xi'+\eta\eta'+\zeta\zeta')\mathord{.}\frac{\mathrm{d}v}{\mathrm{d}\varphi}\\
&=\cos\lambda\lambda'\mathord{.}\frac{\mathrm{d}v}{\mathrm{d}\varphi}
\end{aligned}
\]
Die linke Seite dieser Gleichung, die ebenfalls eine Funktion von \(r,\varphi\) ist, bezeichnen wir mit \(S\). Differenziert man sie nach \(r\), so ergibt sich:
\[
\begin{aligned}
\frac{\mathrm{d}S}{\mathrm{d}r}&=\frac{\mathrm{d}\mathrm{d}x}{\mathrm{d}r^2}\mathord{.}\frac{\mathrm{d}x}{\mathrm{d}\varphi}+\frac{\mathrm{d}\mathrm{d}y}{\mathrm{d}r^2}\mathord{.}\frac{\mathrm{d}y}{\mathrm{d}\varphi}+\frac{\mathrm{d}\mathrm{d}z}{\mathrm{d}r^2}\mathord{.}\frac{\mathrm{d}z}{\mathrm{d}\varphi}\\
&\quad+\tfrac{1}{2}\mathord{.}\frac{\mathrm{d}((\tfrac{\mathrm{d}x}{\mathrm{d}r})^2+(\tfrac{\mathrm{d}y}{\mathrm{d}r})^2+(\tfrac{\mathrm{d}z}{\mathrm{d}r})^2)}{\mathrm{d}\varphi}\\
&=\frac{\mathrm{d}\xi}{\mathrm{d}r}\mathord{.}\frac{\mathrm{d}x}{\mathrm{d}\varphi}+\frac{\mathrm{d}\eta}{\mathrm{d}r}\mathord{.}\frac{\mathrm{d}y}{\mathrm{d}\varphi}+\frac{\mathrm{d}\zeta}{\mathrm{d}r}\mathord{.}\frac{\mathrm{d}z}{\mathrm{d}\varphi}\\
&\quad+\tfrac{1}{2}\mathord{.}\frac{\mathrm{d}(\xi\xi+\eta\eta+\zeta\zeta)}{\mathrm{d}\varphi}
\end{aligned}
\]
Nun gilt \(\xi\xi+\eta\eta+\zeta\zeta=1\), das Differential dieser Summe ist also \(=0\). Wenn \(\rho\) auch hier den Krümmungsradius auf der Linie \(r\) bezeichnet, ergibt sich aus dem vorigen Artikel
\[
\frac{\mathrm{d}\xi}{\mathrm{d}r}=\frac{X}{\rho},\quad\frac{\mathrm{d}\eta}{\mathrm{d}r}=\frac{Y}{\rho},\quad\frac{\mathrm{d}\zeta}{\mathrm{d}r}=\frac{Z}{\rho}
\]
Damit erhalten wir
\[
\frac{\mathrm{d}S}{\mathrm{d}r}=\frac{1}{\rho}\mathord{.}(X\xi'+Y\eta'+Z\zeta')\mathord{.}\frac{\mathrm{d}v}{\mathrm{d}\varphi}=\frac{1}{\rho}\mathord{.}\cos L\lambda'\mathord{.}\frac{\mathrm{d}v}{\mathrm{d}\varphi}=0
\]
denn \(\lambda'\) liegt offenbar auf dem Großkreis mit dem Pol \(L\). Daraus schließen wir, dass \(S\) von \(r\) unabhängig und somit eine Funktion allein von \(\varphi\) ist. Für \(r=0\) gilt aber offenbar \(v=0\) und folglich auch \(\frac{\mathrm{d}v}{\mathrm{d}\varphi}=0\) sowie \(S=0\), unabhängig von \(\varphi\). Daher muss allgemein \(S=0\) gelten, also \(\cos\lambda\lambda'=0\), das heißt \(\lambda\lambda'=90^0\). Daraus schließen wir
